\documentclass[11pt]{article}

\usepackage[margin=0.85in]{geometry}
\usepackage[T1]{fontenc}
\usepackage{amsmath,amssymb,amsthm}
\usepackage{booktabs}
\usepackage{longtable}
\usepackage{array}
\usepackage{enumitem}
\usepackage{hyperref}
\usepackage{microtype}
\usepackage{xcolor}

\emergencystretch=3em
\Urlmuskip=0mu plus 2mu

\hypersetup{
  colorlinks=true,
  linkcolor=blue!50!black,
  urlcolor=blue!50!black,
  citecolor=blue!50!black
}

\setlist[itemize]{leftmargin=1.25em}
\setlist[enumerate]{leftmargin=1.5em}
\newcolumntype{L}[1]{>{\raggedright\arraybackslash}p{#1}}

\newtheorem{law}{Operational law}
\newtheorem{proposition}{Empirical proposition}
\newenvironment{proofsketch}{\paragraph{Evidence sketch.}}{}

\title{\vspace{-2em}What Matters Becomes Measurable:\\
A Synthesis of Structure-Compatible Generalization, Concern Geometry, and Proxy-Resistant Finite Agents}
\author{Jawaun Brown\\Research Derived Experiments}
\date{July 6, 2026}

\begin{document}
\maketitle

\begin{abstract}
This synthesis condenses the July 2026 ``Metaphysics of Intelligence'' archive
and the associated Research Derived Experiments repository into one
reviewer-facing argument. The central result is not that the work proves
consciousness, general agency, or universal OOD robustness. It does something
more actionable: it shows how to test whether finite learning systems succeed
because the right structure controls behavior, rather than because a proxy
passed the surface metric.

Across the archive, the strongest evidence falls into five clusters. First,
structure-compatible model selection shows that, under underspecification,
compatibility with deployment-relevant transformations can select better OOD
models than train or in-distribution validation, with the strongest semantic
selection result reaching selected OOD accuracy 0.978 versus 0.919 for ID,
train, and random selectors. Second, concern geometry shows that value and
viability deform representation: systems allocate resolution, objecthood, and
memory toward what future action requires. Third, the minimal-agent line shows
which mechanisms make concern load-bearing: action-conditioned consequence
models, vector-valued viability heads, identifying interventions, current-error
calibration, value-of-information probes, and decision-layer habituation.
Fourth, Suite C and the long-horizon moved-bottleneck/tool suite turn agency
claims into proxy-resistant benchmarks: behavior alone is not a pass; a pass
requires behavior, a structure-specific gate, and anti-cheat controls. Fifth,
foundation-model contact is bounded but real: open-model and frozen-encoder
results validate selected transport signals, while strict activation-steering
gates reject broader semantic-control claims.

The practical takeaway is a design and evaluation rule: for every proposed
agent capability, find the proxy that can pass without the intended structure,
then add the smallest gate, architecture, or intervention that makes the
intended structure load-bearing. This paper prioritizes the evidence, states
the strongest safe claims, clusters the experiments, gives operational laws
with proof sketches, maps field contributions, identifies the most actionable
and commercially valuable directions, and ranks future work.
\end{abstract}

\section{What Matters First}

Do not read this archive linearly. It contains a research program, a lab
notebook, multiple draft papers, negative controls, reviewer-risk audits, and
publication packages. The fastest accurate read is:

\begin{enumerate}
  \item \textbf{Core technical contribution:} structure-compatible model
  selection under underspecification. When train and ID validation cannot tell
  shortcut and transportable solutions apart, select the model whose learned
  function preserves the transformations deployment will reuse.
  \item \textbf{Core benchmark contribution:} no finite-agent suite should pass
  on behavior alone. Require behavior plus a structure gate plus anti-cheat
  controls.
  \item \textbf{Core architecture contribution:} concern is not a scalar reward
  label. In finite agents it appears as metric allocation, action-conditioned
  consequence prediction, vector viability, identifying intervention,
  re-engagement after world change, and commitment at tool/action surfaces.
  \item \textbf{Core product contribution:} the work is most immediately
  valuable as an eval and observability layer for model selection, agent
  reliability, tool commitment, stale-confidence detection, and OOD governance.
  \item \textbf{Core claim boundary:} the work does not prove consciousness,
  broad autonomy, production reliability, biological validity, or universal OOD
  certification. Its strongest claims are finite, controlled, source-backed,
  and useful.
\end{enumerate}

\paragraph{The one-sentence thesis.}
Finite intelligence becomes measurable when we ask which structure is
preserved, which distinctions the system spends capacity on, and whether that
structure still controls behavior under shift, intervention, memory pressure,
inquiry pressure, and proxy temptation.

\paragraph{The one-sentence evaluation rule.}
A system does not pass because it gets the answer right; it passes when the
right variable is represented, attributed, queried, preserved, or committed at
the causal surface where action is selected.

\section{Source Coverage and Claim Discipline}

The synthesis is based on the local Metaphysics of Intelligence archive and the
repository sources in this worktree. The archive includes the numbered paper
sequence from the geometric-meaning seed paper through the July 6 structure,
benchmark, Suite C, and open-model packages; the ICML publication package; the
Structure-Compatible Generalization folder; Suite C re-engagement and neural
transfer folders; Suite C teacher-free inquiry results; Phase Arc 2 and 3
papers; virtual-governor stress-signal notes; and critical reviews.

Where duplicate PDFs and Markdown/TeX sources exist, this paper treats the
source-backed repository file and the archived PDF as one evidence item. Where
later July 6 versions supersede earlier drafts, the later version supplies the
claim boundary and the earlier versions supply lineage. Negative reports and
reviewer-risk audits are first-class evidence.

\paragraph{Claim levels.}
The archive contains several kinds of evidence:

\begin{itemize}
  \item \textbf{Strong finite empirical result:} preregistered or gated
  controlled experiments with negative controls and result rows.
  \item \textbf{Bounded real-model validation:} open-LM, frozen-encoder, or API
  behavior results that preserve a finite diagnostic.
  \item \textbf{Diagnostic negative:} failed gates that identify a mechanism
  boundary or reject an attractive interpretation.
  \item \textbf{Synthesis or position result:} architecture laws and field
  framing derived from the empirical stack.
  \item \textbf{Conceptual seed:} philosophical vocabulary that motivated the
  tests but is not by itself empirical evidence.
\end{itemize}

The strongest paper should lead with the first two categories, not with the
conceptual seed.

\section{Cluster Map}

\begin{longtable}{L{0.16\linewidth}L{0.29\linewidth}L{0.25\linewidth}L{0.20\linewidth}}
\toprule
\textbf{Cluster} & \textbf{Question} & \textbf{Best evidence} & \textbf{Current claim} \\
\midrule
\endhead
Structure-compatible generalization & Which model will survive deployment
shift when ID evidence ties? & Weakness/OOD papers; SCG phases 1--6; semantic
retrieval and semantic selection; tie-break stress and bootstrap reports. &
Strong finite model-selection claim. \\
\addlinespace
Concern as metric allocation & Does value/viability bend representational
geometry? & Objects from concern; valence tapestry; metric deformation;
effective-dimension law; semantic metric deformation. & Strong controlled
claim, bounded real-encoder validation. \\
\addlinespace
Concern to action & When does representation become competence? & Delta-E
auxiliary learning; two bottlenecks; planning from concern; exploration
diagnostics; state-dependent and allostatic results. & Strong mechanism ladder
with important negative controls. \\
\addlinespace
Self/world attribution and inquiry & How does an agent know what came from
itself, the world, or stale evidence? & First-order self negative; null
intervention; current replay; vector self; world responds; Suite C hand,
learned, teacher-free, source-estimate/tool-transfer results. & Strong finite
agency diagnostic; open-agent replication still needed. \\
\addlinespace
Memory and commitment surfaces & Does a future-critical variable reach the
place where action is committed? & Future Control Moves Memory; Suite D/E tool
commitment; prompt hidden localization; fixed-action localization; causal
patching; API behavior. & Strong synthetic and bounded open-model/API
diagnostic. \\
\addlinespace
Proxy-resistant benchmark methodology & How should agent evals avoid being
fooled by final answers? & Causally grounded agents benchmark; release schema;
benchmark cards; proxy-resistant package. & Highly actionable benchmark
methodology. \\
\addlinespace
Foundation-model boundary & Can hidden semantic geometry steer
generation or action? & Phase 5/6 transport; semantic-specificity tracker;
strict generation/readout gates; Pythia-160M negative replication. & Bounded
transport positives; broad semantic steering not established. \\
\bottomrule
\end{longtable}

\section{Evidence Cluster 1: Structure-Compatible Generalization}

\subsection{What was learned}

The OOD/generalization line began with an operational version of weakness:
measure how many deployment-relevant transformations a candidate function
preserves. The early symbolic result is sharp: in finite cyclic and dihedral
settings where shortcuts and invariant rules both fit training data, weakness
selects the invariant rule in 100\% of trials while the tested simplicity,
description length, MDL-style, compression, validation, loss, and flatness
proxies select the shortcut. Neural sweeps then show learned-function weakness
as a strong OOD correlate.

The newer Structure-Compatible Generalization package broadens this into a
model-selection program. The key shift is from ``weakness as a diagnostic'' to
``compatibility as a selector.'' Across symbolic, modular, vision, learned
generator, finite text, semantic retrieval, and semantic model-zoo selection
settings, the useful question is not ``which model validates best?'' but
``which model preserves the transformation family deployment will require?''

\subsection{Most important results}

\begin{itemize}
  \item Phase 1: true compatibility is the top OOD predictor in modular neural,
  symbolic cyclic, and vision rotation tasks. In selection without OOD labels,
  true compatibility selects mean OOD 1.000 across three domains, versus 0.660
  for train or ID validation and 0.458 for wrong compatibility.
  \item Phase 2: discovered compatibility regularization improves high-ID mean
  OOD from 0.178 to 0.519 and 0.573 at regularization weights 0.05 and 0.20.
  \item Phase 4 language-template substitution: discovered compatibility
  correlates with OOD at Pearson \(r=0.858\), while wrong compatibility is
  weak/negative.
  \item Phase 5 semantic retrieval: true compatibility reaches \(r=0.940\),
  discovered compatibility reaches \(r=0.861\), and wrong compatibility is
  strongly anti-predictive at \(r=-0.872\).
  \item Phase 6 semantic model selection: learned compatibility selects OOD
  0.978, versus 0.919 for random, train, and ID selectors and 0.751 for wrong
  compatibility. The learned-minus-random lift is 0.059 with 95\% CI
  [0.052, 0.065]. Mean-tie, worst-tie, and random-tie stress tests all pass.
\end{itemize}

\subsection{Why it matters}

This is the archive's cleanest main-conference-style contribution. It addresses
underspecification directly: modern pipelines often produce many models with
similar train and validation metrics but different deployment behavior.
Compatibility gives a practical criterion for choosing among them when the
deployment shift has known, inferable, or learnable structure.

\subsection{Boundary}

This is not universal OOD certification. It is a controlled finite result for
structured-domain shifts where the relevant transformation family can be
specified, weakly inferred, or learned well enough to audit. The vision/random
augmentation caveat and phases 1--5 row-regeneration caveats should remain
visible.

\section{Evidence Cluster 2: Concern as Metric Allocation}

\subsection{What was learned}

The conceptual seed says meaning is geometry under concern. The empirical
version is narrower: finite systems spend representational capacity on
distinctions that matter for viability, prediction, or future action.

The object-formation and valence papers show that encoders trained under
viability pressure cluster by causal role rather than sensory similarity.
Vector-valued valence then shows why scalar reward is not enough: preserving
energy and damage as separate predicted consequences supports zero-shot
reweighting across hungry, injured, and balanced contexts, while scalar drive
collapses distinctions that later matter.

The metric-deformation papers extend the same idea into spatial and semantic
representations. Value-weighted training moves local metric density toward
high-value regions. The effective-dimension paper is especially useful because
it is an honest falsification/repair: the expected physical 2-D allocation law
does not hold; the measured exponent implies an effective allocation dimension
near one in the tested RNN harness.

\subsection{Most important results}

\begin{itemize}
  \item Valence encoders group objects by causal reward role, not mere sensory
  color or label.
  \item Vector Delta-V heads preserve multi-dimensional concern and support
  zero-shot priority reweighting where scalar drive heads fail.
  \item Value-weighted metric deformation moves resolution toward priority
  fields in controlled models.
  \item The rate-distortion/effective-dimension test rejects a too-simple 2-D
  allocation law and reveals the architecture's effective bottleneck.
  \item Phase 4 and Phase 6 semantic metric tests show bounded transfer into
  controlled semantic embeddings and actual frozen encoders, with deformed
  margin lift 0.409 versus random-label lift 0.008 in Phase 6.
\end{itemize}

\subsection{Why it matters}

This cluster supplies the scientific bridge between representation learning,
interpretability, and cognitive science. It makes ``what matters'' measurable
as resolution allocation, object formation, value-margin movement, and
priority-sensitive reweighting. For product development, this is the basis of
priority-aware embeddings and retrieval systems: measure whether the embedding
space allocates resolution to the distinctions the task will actually reuse.

\section{Evidence Cluster 3: From Concern to Action}

\subsection{What was learned}

Representation is not competence. Several papers exist because the program
kept finding separations:

\begin{itemize}
  \item a representation can encode reward structure while a policy-gradient
  learner fails to exploit it;
  \item a policy can succeed while using a proxy representation;
  \item a model can predict total outcome while misattributing self and world
  components;
  \item an uncertainty signal can be high while not indicating useful
  information gain;
  \item a quiet probe policy can be stable or falsely calm.
\end{itemize}

The clean positive mechanism is \textbf{predictive policy closure}: an
action-conditioned consequence model can become the policy by argmax. In
Planning from Concern, a Delta-E model reaches full return and near-perfect
action accuracy on the XOR task without optimal-action labels or
policy-gradient training.

The negative results are just as important. State-dependent concern fails under
smooth online training at the singular boundary. Off-policy state coverage
helps but does not solve the boundary. An oracle boundary feature or learned
hard gate solves the task. Ensemble uncertainty fails when all models share the
same blind spot. Allostatic regulation helps most when it is simple and greedy,
not when uncertainty bonuses amplify wrong confidence.

\subsection{Why it matters}

This cluster contributes practical architecture rules. Learn consequences, not
only actions. Preserve vector-valued consequences until decision time. Ensure
action coverage. Represent hard regime variables explicitly. Treat ensemble
variance as a hypothesis, not as epistemic truth.

\section{Evidence Cluster 4: Self/World Attribution and Adaptive Inquiry}

\subsection{What was learned}

The self/world line is one of the archive's most important scientific arcs.
The first result is negative: architectural factorization alone does not
recover a self/world decomposition. A self head that sees actions and a world
head that does not can still split total prediction arbitrarily while preserving
behavior. This is the gauge-symmetry problem.

Null interventions break that gauge. A null action that is dynamically like
skip but labeled as a world-anchor reduces false credit dramatically. But the
next problem is deciding when to spend costly null probes. Early V-probe
targets fail by saturating, tracking stale residuals, or confusing shock noise
with systematic error. Current-replay calibration fixes part of the problem by
recomputing residuals against the current model. Scale normalization fixes a
vector-dimensional calibration failure.

When the world responds to the agent, the decomposition must become richer:
direct self, mediated world, and exogenous world. Current error stops being a
good proxy for probe value. The agent needs value-of-information. This leads to
the Suite C line: re-engage after world change, recover attribution, avoid
false calm, spend probes selectively, and reopen after a second shift.

\subsection{Most important results}

\begin{itemize}
  \item First-order self: factorized heads fail; behavior can be correct while
  attribution is false.
  \item Null intervention: active world anchoring reduces false credit; passive
  exposure is insufficient.
  \item Current-error calibration: current replay reduces food self MAE close
  to oracle and beats matched random probing.
  \item World Responds: three-head direct/mediated/exogenous decomposition is
  the right structural choice when the world is action-correlated.
  \item Suite C hand policy: final affected MAE 0.112, first-shift selectivity
  16.927, second-shift reopenability 16.667, 22.6 probes, all C1--C6 gates pass.
  \item Suite C learned transfer: final affected MAE 0.112, first-shift
  selectivity 16.667, second-shift reopenability 17.448, 23.1 probes; stale,
  wrong, suppressed, and matched-random controls fail.
  \item Suite C teacher-free: reward/CEM policy passes C1--C6 plus T1/N1 with
  final affected MAE 0.095, recovery 1.000, selectivity 12.500, second-shift
  reopenability about 11, and about 22 probes. A 64-seed bootstrap preserves
  the pass, with selectivity lift over matched random 11.292 [10.976, 11.596].
  \item Source-estimate/tool-transfer follow-up: replacing privileged
  source-identity with an observable estimate and routing through a local
  JSON-like tool adapter passes, while malformed-tool control fails.
\end{itemize}

\subsection{Why it matters}

This is the archive's best agency-eval result. It gives an operational test for
world-coupled agency: not just acting, not just remembering, but knowing when
the model has lost contact with the world and reopening inquiry selectively.
For product development, this points toward a \textbf{re-engagement governor}
between memory, tools, planning, and action.

\section{Evidence Cluster 5: Memory and Commitment Surfaces}

\subsection{What was learned}

Future Control Moves Memory introduces the long-horizon moved-bottleneck
diagnostic. Four early slots are matched for salience and frequency. One
becomes future-critical, and the critical slot moves across runs. The question
is not only whether the agent answers correctly. It is whether hidden-state
sensitivity, generated action, tool call, repair behavior, and causal readout
move with the future-critical slot.

The result ladder is unusually complete. The diagnostic passes through a
synthetic trained-agent ladder, longer horizons, closed-loop tool handoff,
repair, structured tool calls, multi-field schemas, stochastic API failures,
8-slot arguments, aliases, text arguments, generated JSON, autoregressive JSON,
prompt-level open models, hidden localization, fixed-action localization,
causal patching, prompt-family robustness, and black-box API behavior.

\subsection{Most important results}

\begin{itemize}
  \item Synthetic trained agents reach 1.000 closed-loop accuracy and positive
  moved-slot specificity across multiple tool/repair/action surfaces.
  \item Initial prompt-level final-prompt hidden specificity is a controlled
  strong negative: behavior passes, but the measured hidden site does not.
  \item Hidden localization resolves the negative: 17 registered hidden sites
  pass across Qwen 0.5B, Qwen 1.5B, and SmolLM2 1.7B, strongest at generated
  final states.
  \item Fixed-action localization removes the generated-answer-token confound:
  24 hidden sites pass even when assistant action tokens are held fixed.
  \item Fixed-prefix causal patching shows value-prefix hidden states can shift
  next-token value logits, while prompt-final sites remain negative controls.
  \item Prompt-family causal patch robustness passes across standard, compact,
  and audit-checklist framings.
  \item API black-box behavior: Gemini, Anthropic, and OpenAI pass the matched
  prompt-family suite. External stress is mixed: Gemini and Anthropic pass,
  OpenAI GPT-4.1 Nano shows a sparse dispatch failure localized to one
  reproduced cell out of 16 in the robustness follow-up.
\end{itemize}

\subsection{Why it matters}

This cluster gives an interpretable and product-relevant memory test. A
business agent should not only hold context; it should preserve the variable
that will later control a tool call, schema field, repair branch, or emitted
value. The long-horizon suite is directly actionable for tool-use reliability,
workflow agents, API-call auditing, and context-window memory diagnostics.

\section{Evidence Cluster 6: Foundation-Model Contact and Rejected Steering}

\subsection{What was learned}

The archive contains a tempting story about hidden semantic geometry in
language models. The careful version is: hidden geometry often predicts labels
or moves log-prob surfaces, but strict behavior, generation, alias, and control
gates reject broad semantic steering claims.

Phase 5 is a transport suite. It shows that controlled proxy mechanisms survive
harder model-like gates: language-action transport passes in a stronger proxy,
semantic metric deformation survives a foundation-style encoder proxy,
role-routed heads break the mediated-identifiability ceiling, and seam
consistency, not topology alone, carries OOD lift.

Phase 6 moves to actual open LMs and frozen encoders. It reports weak positive
open-LM action-coupling signals across three decoder LMs, with Qwen 0.5B
strongest and Pythia 70M weak/control-like. Frozen-encoder metric deformation
passes with deformed margin lift 0.409 versus random 0.008.

The activation-steering line is different. The semantic-specificity tracker
records many failed rescues: held-out aliases, target-disjoint controls,
random relation nulls, generation-match gates, learned generation readout,
yes-bias controls, residualization, whitening, pair-specific optimization,
feature masks, state gates, multi-class prototype gates, and Pythia-160M
replication all prevent a broad steering claim. The current safe statement is
diagnostic: the strict verifier exposes a narrow specificity frontier, but the
frontier has not shown stable held-out alias specificity on expanded controls.

\subsection{Why it matters}

This is exactly how a healthy research program should behave. It promotes
bounded transport signals while rejecting overbroad mechanistic claims. For
future research, exact-relation readout identifiability should come before more
steering optimization on the same brittle interface.

\section{Operational Laws and Proof Sketches}

\begin{law}[Behavior is not load-bearing structure]
Final accuracy, return, or tool success is not sufficient evidence that a
system used the intended structure.
\end{law}
\begin{proofsketch}
The archive repeatedly constructs behavior-only passes that fail structure
gates: shortcut OOD rules, gauge-symmetric self/world heads, visible-control
memory cases, signal-suppression false calm, and activation directions that
move positives and controls together. The proxy-resistant benchmark rule is the
repair: behavior plus structure gate plus anti-cheat control.
\end{proofsketch}

\begin{proposition}[Structure-compatible selection]
In finite structured shifts, compatibility with the deployment-relevant
transformation family can improve OOD-free model selection among train/ID-tied
candidates.
\end{proposition}
\begin{proofsketch}
The evidence ladder runs from oracle finite transformations through inferred
and learned generators, finite text substitutions, frozen semantic retrieval,
and semantic model-zoo selection. Phase 6 selects OOD 0.978 with discovered
compatibility versus 0.919 for train/ID/random and 0.751 for wrong
compatibility, with bootstrap and tie-break stress passing. The boundary is
finite structured shifts, not universal OOD certification.
\end{proofsketch}

\begin{law}[Predictive policy closure]
When an agent learns an action-conditioned consequence model for what it cares
about, the consequence model can become the policy by action argmax.
\end{law}
\begin{proofsketch}
Planning from Concern shows that a Delta-E consequence model solves the XOR
homeostatic task with near-perfect action accuracy and full return, without
optimal-action labels or policy-gradient training. Earlier Delta-E auxiliary
and two-bottleneck papers show why this matters: representation alone is not
competence, and policy-gradient learning can fail to exploit a useful
representation.
\end{proofsketch}

\begin{law}[Vector concern must survive until decision time]
Scalar reward collapses distinctions that later priorities may need.
\end{law}
\begin{proofsketch}
The valence tapestry and vector first-order self papers show that vector
Delta-V heads support zero-shot reweighting across hungry, injured, and
balanced contexts, especially for medicine handling. Scalar drive heads cannot
preserve the needed distinctions.
\end{proofsketch}

\begin{law}[Attribution needs identifying interventions]
Architectural factorization is not self/world identification.
\end{law}
\begin{proofsketch}
First-Order Self shows that self/world heads can split total prediction
arbitrarily while preserving behavior. Null interventions, source labels,
temporal asymmetry, or interventional contrast are needed to break the gauge.
Null anchoring and later contrastive/role-routed mechanisms supply positive
evidence.
\end{proofsketch}

\begin{law}[Probe value is not current error]
The best next question is not necessarily where the current model is most
wrong.
\end{law}
\begin{proofsketch}
World Responds shows an oracle-current-error probe can be much worse than a
learned probe because high error may be irreducible or poorly positioned for
future model improvement. Phase 4's learned VOI gate passes while current error
and ensemble variance fail. Suite C strengthens this into re-engagement under
world change.
\end{proofsketch}

\begin{law}[Cool the decision, not the signal]
Habituation should reduce unnecessary action while preserving surprise.
\end{law}
\begin{proofsketch}
Habituated Re-engagement shows decision-layer cooling reduces anxious probing,
while signal-layer cooling creates false calm. Suite C formalizes the no-false
calm gate: probe rate, uncertainty signal, and outcome error must improve
together.
\end{proofsketch}

\begin{law}[Memory matters at commitment surfaces]
Memory becomes agent-relevant when a future-critical variable reaches an
action, tool, schema, repair, emitted value, or causal readout surface.
\end{law}
\begin{proofsketch}
The moved-bottleneck suite shows memory sensitivity follows the moved critical
slot through hidden states, generated JSON, tool calls, repair branches, and
causal value-token patches. Visible-control and prompt-final negative controls
rule out several shortcut explanations.
\end{proofsketch}

\section{Field Contributions}

\subsection{Machine learning generalization}

The contribution is a model-selection criterion for underspecification. If
many models tie on train and ID validation, select by compatibility with the
deployment-generating structure. This directly targets shortcut learning,
domain generalization, and OOD model selection.

\subsection{Agent evaluation and AI safety}

The benchmark contribution is the minimum pass rule. A capable-looking agent
should not pass unless the right causal/structural reason explains its behavior.
This connects to reward hacking, goal misgeneralization, stale confidence,
tool-call brittleness, and proxy optimization.

\subsection{Interpretability}

The work shifts interpretability from ``is information present?'' to ``does the
information control behavior at the right surface?'' The memory, attribution,
causal patch, and activation-specificity lines are all examples of this
surface-based standard.

\subsection{Agent architecture}

The archive suggests concrete modules: action-conditioned consequence models,
vector viability heads, source-attribution gates, null or contrastive
interventions, value-of-information probe policies, re-engagement governors,
role-routed heads, and commitment-surface tests.

\subsection{Cognitive science, neuroscience, and philosophy}

The contribution is not a proof of consciousness. It is a computational
operationalization of concern-like organization: what matters bends
representation, inquiry, action, attribution, and memory. This gives philosophy
and cognitive science a testable vocabulary while keeping biological and human
claims out of scope.

\section{Actionability, Ironcladness, and Commercial Value}

\begin{longtable}{L{0.23\linewidth}L{0.22\linewidth}L{0.22\linewidth}L{0.23\linewidth}}
\toprule
\textbf{Area} & \textbf{How ironclad?} & \textbf{How actionable?} & \textbf{Why lucrative?} \\
\midrule
\endhead
Structure-compatible model selection & High within finite structured shifts;
Phase 6 has bootstrap and tie-break stress. & Add compatibility selectors to
model-zoo and eval pipelines. & Better model selection under deployment risk
is immediately valuable for regulated and enterprise ML. \\
\addlinespace
Proxy-resistant agent benchmarks & High as methodology; suites vary by
maturity. & Ship benchmark cards, JSONL rows, and one-command smoke tests. &
Every agent platform needs evals that catch shortcut success and stale
confidence. \\
\addlinespace
Suite C re-engagement governor & High in finite harness; open-agent transfer
pending. & Add stale-belief monitors and selective inquiry policies between
memory and action. & Reduces costly agent mistakes from stale assumptions,
especially in long-running workflows. \\
\addlinespace
Commitment-surface memory tests & High across synthetic, open-model, and API
behavior surfaces, with clear non-claims. & Test whether critical variables
survive tool calls, schemas, repairs, and generated values. & Tool-use and
workflow reliability are direct product pain points. \\
\addlinespace
Semantic metric deformation & Medium-high; controlled and frozen-encoder
evidence is positive. & Build priority-aware retrieval or embedding audits. &
Retrieval quality and high-value-case recall are commercially important. \\
\addlinespace
Activation steering & Low for broad claims; high as a diagnostic negative. &
Stop optimizing brittle global vectors; test exact readout first. & Avoids
shipping unsafe or fake controls; useful for research tooling, not product yet. \\
\addlinespace
Concern-geometry synthesis & Medium as theory; depends on narrower papers. &
Use as architecture vocabulary after empirical anchors. & Valuable for
strategy, but not a standalone product proof. \\
\bottomrule
\end{longtable}

\section{Product Development Implications}

\paragraph{1. Model selection product.}
Build an OOD-free selector that evaluates compatibility against known or
inferred deployment transformations. Inputs: candidate models, deployment
transform hypotheses, train/ID filters, wrong-transform controls. Output:
selected model, regret estimate, compatibility audit, claim boundary.

\paragraph{2. Agent reliability product.}
Build a proxy-resistant eval suite with vector scores: behavior, structure,
attribution, inquiry, commitment, and generalization. The minimum pass rule
should be a product primitive, not a paper footnote.

\paragraph{3. Re-engagement governor.}
Add a layer between memory, tools, and action that tracks assumption staleness,
world-change signals, affected-source estimates, cost-aware inquiry, and
no-false-calm recovery. The first product version can be benchmark-only; the
second can become runtime monitoring.

\paragraph{4. Tool commitment observability.}
Instrument agents to verify that future-critical variables survive prompt
state, tool-call JSON, schema validation, repair branches, and emitted values.
The long-horizon suite is a direct template.

\paragraph{5. Priority-aware retrieval.}
Use metric-deformation audits to determine whether retrieval/embedding systems
allocate resolution to high-value business distinctions, not only average
semantic similarity.

\section{Prioritized Future Directions}

\begin{enumerate}
  \item \textbf{External Suite C replication.} Highest scientific and product
  priority. Reproduce teacher-free/source-estimate/tool-transfer Suite C in an
  open LM or API-agent harness. This would move the inquiry governor from a
  finite simulator result to a broadly legible agent-eval result.
  \item \textbf{Publish Structure-Compatible Model Selection.} Keep Phase 6 as
  the headline, include row-level artifacts, bootstrap intervals, tie-break
  stress, formal compatibility definitions, and a strict finite-structured
  claim boundary.
  \item \textbf{Release proxy-resistant benchmark package.} Package Suite C and
  Suite D/E with public JSONL rows, benchmark cards, fixture smoke tests,
  release schema, and a Hugging Face dataset-style card.
  \item \textbf{Productize commitment-surface memory tests.} Turn the
  moved-bottleneck/tool suite into a reusable CI/eval for tool-using agents.
  \item \textbf{Build role-routed self/world architectures.} Phase 4 suggests
  the mediated-identifiability ceiling is architectural. Test disjoint
  per-role heads, mixture-of-experts, and richer counterfactual interventions
  in larger agent environments.
  \item \textbf{Turn compatibility into training, not only selection.} Extend
  discovered compatibility regularization to practical domains where
  deployment transforms are partially known.
  \item \textbf{Advance priority-aware embedding audits.} Move semantic metric
  deformation from post-hoc/frozen evaluation toward controlled finetuning or
  adapter training with business-relevant priority fields.
  \item \textbf{Do exact-relation readout before more steering.} The activation
  line should first establish stable exact-relation identifiability under
  strict controls before attempting stronger generation steering.
  \item \textbf{Keep the broad concern-geometry paper as synthesis.} It should
  follow the narrower empirical submissions. It is valuable, but leading with
  it risks reviewer confusion.
\end{enumerate}

\section{What Not To Prioritize}

\begin{itemize}
  \item Do not lead with consciousness, AGI, or metaphysics in technical
  venues.
  \item Do not claim universal OOD certification.
  \item Do not claim production agent reliability from finite Suite C or
  black-box API behavior alone.
  \item Do not promote activation steering as a positive result yet.
  \item Do not collapse benchmark results into a single leaderboard score.
  \item Do not hide negative controls; they are the reason the program is
  credible.
\end{itemize}

\section{Recommended Paper and Release Sequence}

\begin{enumerate}
  \item \textbf{Structure-Compatible Model Selection Under Underspecification.}
  Main-conference target. Most complete, most legible, strongest statistical
  anchor.
  \item \textbf{Proxy-Resistant Benchmarks for Causally Grounded Finite
  Agents.} Benchmark target. Lead with the minimum pass rule and Suite C plus
  Suite D/E.
  \item \textbf{Teacher-Free Re-Engagement.} Focused agency paper. Stronger if
  external open-agent replication is added.
  \item \textbf{What Matters Becomes Measurable.} Broad synthesis/position
  paper. Use after the narrower empirical anchors exist publicly.
\end{enumerate}

\section{Conclusion}

The archive's deepest contribution is not any single experiment. It is a
research discipline: whenever a system appears capable, ask what proxy could
make the behavior pass without the intended structure. Then make the structure
load-bearing, add an anti-cheat control, and bound the claim to the weakest
validated level.

That discipline produced several strong results. Structure-compatible model
selection gives a practical answer to underspecification. Suite C gives a
finite test of adaptive inquiry after world change. The moved-bottleneck suite
shows how future-control relevance moves memory and tool commitment. The
metric-stack line shows how concern deforms representation and where
architecture ceilings appear. The activation-steering failures show that the
program can reject an attractive story when the controls do not support it.

If one sentence should survive into future research and product work, it is
this: \textbf{Do not ask only whether the agent succeeded; ask whether the
future-controlling structure reached the surface where success was chosen.}

\appendix

\section{Primary Reading Spine}

The shortest serious reading order is:

\begin{enumerate}
  \item ICML publication package README.
  \item Structure-compatible model-selection result tables.
  \item Semantic-selection control paper.
  \item Causally grounded finite-agent benchmark paper.
  \item Suite C teacher-free inquiry package README.
  \item Suite C teacher-free wide-seed bootstrap report.
  \item Long-horizon moved-bottleneck paper and benchmark card.
  \item Metric Stack synthesis paper.
  \item Paper-readiness and semantic-specificity audit notes.
\end{enumerate}

\section{Source Clusters Used}

\begin{itemize}
  \item \textbf{Conceptual and metric-stack archive:}
  \url{0_geometric_meaning_and_agency.pdf};
  \url{26_Metric_Stack_of_Concern_v4_2026_07_06.pdf};
  \url{papers/metric_stack_synthesis/paper.md};
  metric-stack position package.
  \item \textbf{Weakness and structure-compatible generalization:}
  \url{1_Weakness_Predicts_OOD_2026_06_09.pdf};
  \url{30_Weakness_Predicts_OOD_2026_07_02_draft.pdf};
  weakness-invariance paper;
  structure-compatible generalization papers, results, and ICML package.
  \item \textbf{Concern/action/mechanism ladder:}
  numbered papers 4--25 in the archive; Planning from Concern; First-Order
  Self; Null Intervention; World Responds; Probe-Value Re-engagement; and
  Habituated Re-engagement.
  \item \textbf{Suite C:}
  Suite C re-engagement archive; neural probe transfer archive; teacher-free
  inquiry package; and world-responds result reports.
  \item \textbf{Long-horizon and tool commitment:}
  \url{31_Future_Control_Moves_Memory_2026_07_06.pdf}; long-horizon
  moved-bottleneck paper; benchmark card; and result reports.
  \item \textbf{Foundation-model and activation boundary:}
  Phase 4, Phase 5, and Phase 6 archive papers; paper-readiness and
  semantic-specificity audits; and activation-geometry result reports.
  \item \textbf{Benchmark packaging and release:}
  Causally grounded agents benchmark paper; benchmark docs; release schema;
  publication sharing map; and proxy-resistant agent benchmark package.
\end{itemize}

\section{Safety and Non-Claims}

This synthesis intentionally avoids the following claims:

\begin{itemize}
  \item consciousness, sentience, or human-like subjective experience;
  \item biological or neural validity beyond computational analogy;
  \item production reliability for deployed agents;
  \item universal OOD certification;
  \item open-world semantic robustness;
  \item broad activation-steering control.
\end{itemize}

\end{document}
